\begin{trapbox}

	\begin{description}[style=nextline, leftmargin=2.8em, labelsep=0.5em, itemsep=0.35em]
		\item[\textbf{\textcolor{CTrap}{易错点一（缺失线在面内）：}}]
			学生误以为只要一条直线垂直于一个平面，就能推出该直线所在的任何平面都与那个平面垂直，即忽略直线必须在另一个平面内这一条件。
		\item[\textbf{\textcolor{CTrap}{正确理解（必要条件）：}}]
			判定定理要求 $l\subset\beta$ 是必不可少的条件。缺少它，结论不一定成立。
		\item[\textbf{\textcolor{CTrap}{错因分析（定理记忆不全）：}}]
			对定理“若 $l\perp\alpha,\ l\subset\beta$，则 $\alpha\perp\beta$”记忆不准确，丢掉“$l\subset\beta$”这一关键条件。
	\end{description}

	\medskip
	\begin{description}[style=nextline, leftmargin=2.8em, labelsep=0.5em, itemsep=0.35em]
		\item[\textbf{\textcolor{CTrap}{易错点二（混淆线面垂直）：}}]
			学生将“直线垂直于平面内的一条直线”误认为线面垂直，从而错误地判定面面垂直。
		\item[\textbf{\textcolor{CTrap}{正确理解（线面垂直定义）：}}]
			线面垂直要求直线垂直于平面内的任意一条直线，等价于垂直于该平面内两条相交直线。
		\item[\textbf{\textcolor{CTrap}{错因分析（判定定理模糊）：}}]
			对线面垂直的判定定理掌握不牢固，只注意到一条垂直就以为得到线面垂直。
	\end{description}

	\medskip
	\begin{description}[style=nextline, leftmargin=2.8em, labelsep=0.5em, itemsep=0.35em]
		\item[\textbf{\textcolor{CTrap}{易错点三（逆转逻辑错误）：}}]
			学生误认为由 $\alpha\perp\beta$ 可以推出 $\beta$ 内任意直线都垂直于 $\alpha$。
		\item[\textbf{\textcolor{CTrap}{正确理解（面面垂直性质）：}}]
			面面垂直的性质是：若 $\alpha\perp\beta$，$\alpha\cap\beta=m$，$l\subset\beta$，$l\perp m$，则 $l\perp\alpha$。不是任意直线都垂直。
		\item[\textbf{\textcolor{CTrap}{错因分析（性质与判定混淆）：}}]
			将面面垂直的判定与性质中的条件混淆，认为只要面面垂直就能得到线面垂直。
	\end{description}

\end{trapbox}
